% SPDX-License-Identifier: LPPL-1.3c
% Copyright (C) 2026 Christophe Jorssen
% Author and Current Maintainer: Christophe Jorssen <christophe.jorssen@gmail.com>
% LPPL maintenance status: maintained
% This file is part of luafemm. See LICENSE and LICENSES.md for its terms.
\section{Nested domains and compound paths}
\label{sec:domains}
Two rectangles inside one another can describe a core and its surrounding
material. There are two complementary ways to express this in TikZ:
paint the outer region first and the inner region second, or give the
surrounding material an explicit hole. Both use the same conforming mesh.

\subsection{Two ordinary paths: later declarations win}
With separate region paths, the physical declaration order matters.
The last region containing a point supplies its material, current density
and magnetisation direction. A broad outer region therefore comes first:
\begin{codeexample}[width=5cm]
\begin{tikzpicture}[femm/problem={
  xmin=-25,xmax=25,ymin=-20,ymax=20,
  mesh size=3}]
  \filldraw[femm/region={
    material=pure-iron},fill=gray!25]
    (-20,-15) rectangle (20,15);
  \filldraw[femm/region={
    material=cast-alnico-5-lng37,
    magnetization angle=90},fill=red!25]
    (-10,-6) rectangle (10,6);
  \pic {femm mesh};
\end{tikzpicture}
\end{codeexample}
The iron occupies the space between the rectangles; alnico occupies the
inner rectangle. Reversing these two declarations would replace the
alnico with iron. Graphical opacity, clipping and drawing layers do not
perform material subtraction. An explicit later |material=air| region
paints air in exactly the same way.

\subsection{A surrounding material with its own hole}
To express the annulus directly, use one physical path with two closed
subpaths and TikZ's native |even odd rule|. No additional luafemm fill-rule
key is needed. The same rule controls the visible fill and the material:
\begin{codeexample}[width=5cm]
\tikzset{iron shell/.style={
  even odd rule,
  femm/region={material=pure-iron},
  draw,fill=gray!25}}
\begin{tikzpicture}[femm/problem={
  xmin=-25,xmax=25,ymin=-20,ymax=20,
  mesh size=3}]
  \filldraw[femm/region={
    material=cast-alnico-5-lng37,
    magnetization angle=90},fill=red!25]
    (-10,-6) rectangle (10,6);
  \path[iron shell]
    (-20,-15) rectangle (20,15)
    (-10,-6) rectangle (10,6);
  \pic {femm mesh};
\end{tikzpicture}
\end{codeexample}
Here the magnet was declared first. It survives because the iron path
does not cover the cavity. The cavity preserves the previously declared
material; if none was declared there, it contains the background air.
A later material may still paint over any part of this compound region.

\subsection{The native TikZ fill rules}
\begin{key}{/tikz/even odd rule}
For a physical region, count how many contour boundaries a ray from the
point crosses. Odd counts fill the region, even counts leave it unchanged.
Contour direction and the order of subpaths within this one path do not
matter. With three nested contours, the outer band and innermost island
receive the material, while the intermediate cavity preserves its previous
material. Disjoint contours describe separate pieces with common properties.
\end{key}
\begin{key}{/tikz/nonzero rule}
PGF's default rule uses a signed winding number. A point is filled when
that number is nonzero. Nested contours with the same direction therefore
fill the centre too; reverse the inner contour to make a hole. For nested
rectangles or circles, |even odd rule| is usually easier to read. A native
fill-rule key on a path overrides an inherited rule in the ordinary TikZ way.
\end{key}
These are existing PGF/TikZ keys, documented here for their additional
physical effect. luafemm reads the evaluated PGF fill rule during capture.
It does not depend on whether the path also requests graphical filling:
|\path[even odd rule,femm/region=...]| has the same material membership
as |\filldraw| with those options.

Every subpath must close, using |rectangle|, |circle|, |ellipse| or
|-- cycle|. Each individual contour must remain simple: self-intersections,
self-contact and backtracking are rejected after flattening. Distinct
contours may intersect or share edges; the constrained mesher subdivides
their intersections before classification. Curves use the existing
|curve tolerance| and coordinate-rounding conventions. Features lost in
that approximation cannot be recovered by refining the mesh.

All filled pieces of one compound region share its material, current
density, magnetisation direction and local mesh size. Total source current
is current density times the final painted area, not the sum of the
unsigned contour areas. Source current is not automatically renormalised
when another region paints over a winding. To use different properties
in several pieces, declare separate physical paths.

\subsection{Meshing, field sampling and cache reuse}
Both meshers include every subcontour as an interface constraint. The
grid mesher still requires axis-aligned straight edges. Delaunay local
refinement applies inside the filled portion and along all its contour
edges, including cavity walls. It does not request the same fine interior
lattice throughout a hole; another region can refine that hole if desired.
Boundary point spacing can nevertheless influence neighbouring triangles.

The full rectangular problem domain remains meshed. A hole in a material
is \emph{not an unmeshed exclusion}: air and any inner media still carry
field, and a measurement profile can cross them. No boundary condition
is added on the cavity wall. The ordinary material-interface transmission
conditions follow from the finite-element formulation. Sampling exactly
on an interface retains the documented first-containing-triangle
convention; use a profile offset when a particular side matters.

Unmeshed holes would change the PDE domain and require their own boundary
conditions and solvability checks. They are not introduced by |even odd rule|
or by a TikZ clip, and remain unsupported in this version. Measurement
profiles also retain a separate contract: each profile is one connected
curve, not a compound region boundary. Give separate contours separate
profile names.

The cache descriptor contains every contour and the fill rule. Changing
either invalidates the mesh, because local refinement also depends on
membership. Material-only changes can still reuse geometry. The cache
revision is now 4, including quality-refinement controls and diagnostics;
|auto| rebuilds earlier records and |frozen| rejects them.
The complete \texttt{examples/nested-domains.tex} example combines an alnico
magnet, an air cavity and a nonlinear iron shell with field lines and a
PGFPlots scan across all three media.

\subsection{Complete example: a magnet inside a shell}
The following source declares the alnico first and the compound iron
shell second. Its hole preserves both the inner magnet and the surrounding
air. The scan crosses the shell, air and magnet; the plotted jumps reflect
their different material laws and the elementwise field approximation.
The numbered comments also make the declaration/solve/plot order explicit.
\luafemmexample{nested-domains}
